% Section content template for individual Educative lessons
% This file contains the actual content for a specific section
% Generated from Educative API section content

% Section: Evaluation of Google Docs’ Design
% Section ID: 5489039691481088
% Section Slug: evaluation-of-google-docs-design
% Generated: 2025-12-08T15:23:31.535583

\subsection{Requirements compliance}\label{0L6s9njl3vCo4TAgnPElM}
\\
This lesson will focus on how our design meets the non-functional requirements. In particular, we'll focus on consistency, latency, scalability, and availability.

\subsubsection{Consistency}\label{flWOWkjaAUxejHeTb9Bjf}
\\
We've looked at how we'll achieve strong consistency for conflict resolution in a document through two technologies: operational transformation (OT) and Conflict-free Resolution Data Types (CRDTs). In addition, a time series database enables us to preserve the order of events. Once OT or CRDT has resolved any conflicts, the final result is saved in the database. This helps us achieve consistency in terms of individual operations.
\\
We're also interested in keeping the document state consistent across different servers in a data center. To replicate an updated state of a document within the same data center at the same time, we can use peer-to-peer protocols like Gossip protocol(Also referred to as epidemic protocol which is used for disseminating data across all nodes in a network to achieve consistency.). Not only will this strategy improve consistency, it will also improve availability.

\begin{quote}
\textbf{Quiz:}
\textbf{Question 1:} Why should we use strong consistency instead of eventual consistency for conflict resolution in a collaborative document editing service?
\textbf{Answer:} From Amazon's Dynamo system, we learn that if we use eventual consistency for conflict resolution, we might have multiple versions of a document that are eventually reconciled, either automatically or manually. In the case of automatic reconciliation, the document might update abruptly, which defeats the purpose of collaboration. The second case, manual resolution, is tedious labor that we want to avoid.

Therefore, we use strong consistency for conflict resolution, and the logically centralized server provides the final order of events to all clients. We use a replicated operations queue so that even if our ordering service dies, it can easily restart on a new server and resume where it left off. Clients might encounter short service unavailability while the failed component is being respawned.
\end{quote}

\subsubsection{Latency}\label{u6wU_1KU44yhtX8PDvZm6}
\\
Latency may feel like a challenge, specifically when two users are distant from each other or the server. However, users maintain a replica of documents at their end while data is being propagated through WebSockets to end-servers. Therefore, user-perceived latency will be low. Apart from this, users mostly tend to write textual data in a small document. Therefore, dissemination of data among different servers within the same facility and among different data centers or zones will be carried out with low latency. Also , files like videos and images can be stored in CDNs for quick serving because this content doesn\textquotesingle t change frequently.\\
\strut \\
Practically speaking, there will be a limited number of readers and writers of an online document. For readers specifically, latency won\textquotesingle t be an issue because the document will be loaded only once. So , most readers can be served from the same data center. For writers, an optimal zone should be selected as the centralized location between collaborators of the same document. However, for popular documents, asynchronous replication will be an effective method to achieve good performance and low latency for a large number of users. In general, achieving strong consistency becomes a challenge when replication is asynchronous.

\subsubsection{Availability}\label{c_F1Yp9kSZNL0igmsOVCl-1}
\\
Our design ensures availability by using replicas and monitoring the primary and replica servers using monitoring services. Key components like the operations queue and data stores internally manage their replication.
\\
Since we use WebSockets, our WebSocket servers can connect users to the session maintenance servers that will determine if a user is actively viewing or collaborating on a document. Therefore, keeping multiple WebSocket servers will increase the availability of the design. Lastly , we employ caching services and CDNs to improve availability in case of failures.
\\
However, at the moment, we haven\textquotesingle t devised a disaster recovery management scheme.

\textbf{Non-functional Requirements Compliance}

\begin{table}[htbp]
\centering
\small
\adjustbox{max width=\textwidth}{
\begin{tabular}{|p{0.195\textwidth}|p{0.755\textwidth}|}
\hline
\textbf{Requirements} & \textbf{Techniques} \\
\hline
\hline
Consistency & \begin{itemize} \tightlist \item Gossip protocol to replicate operations of a document within the same data center \item Concurrency techniques like OT and CRDTs \item Usage of time series database for maintaining the order of operations \item Replication between data centers \end{itemize} \\
\hline
Latency & \begin{itemize} \tightlist \item Employing WebSockets \item Asynchronous replication of data \item Choosing optimal location for document creation and serving \item Using CDNs for serving videos and images \item Using Redis to store different data structures including CRDTs \item Appropriate NoSQL databases for the required functionality \end{itemize} \\
\hline
Availability & \begin{itemize} \tightlist \item Replication of components to avoid SPOFs \item Using multiple WebSocket servers for users that may occasionally disconnect \item Component isolation improves availability \item Implementing disaster recovery protocols like backup, replication to different zones, and global server load balancing \item Using monitoring and configuration services \end{itemize} \\
\hline
Scalability & \begin{itemize} \tightlist \item Different data stores for different purposes enable scalability \item Horizontal sharding of RDBMS \item CDNs capable of handling a large number of requests for big files \end{itemize} \\
\hline
\end{tabular}
}
\caption{Non-functional Requirements Compliance}
\end{table}

\subsubsection{Scalability}\label{c_F1Yp9kSZNL0igmsOVCl-2}
\\
Since we\textquotesingle ve used microservice architecture, we can easily scale each component individually in case the number of requests on the operations queue exceeds its capacity. We can use multiple operations queues. In that case, each operations queue will be responsible for a single document. We can forward operations requested by different users that are associated with a single document to a specific queue. The number of spawned queues will be equal to the number of active documents. As a result, we're able to achieve horizontal scalability.
\\
Imagine you are designing a collaborative document editing tool for a global team. Some users work in areas with high-speed internet, while others face slower or unreliable connections. A consistent user experience means waiting for updates from everyone, but faster responsiveness may lead to temporary inconsistencies between users' views of the document.\\
\strut \\
Should the system prioritize real-time speed or consistency for all users, and why?\\
\textbf{\hfill\break
} Give your reasoning in the AI assessment widget below.

\subsection{Conclusion}\label{jZ3utsWz3CdBm6fuUaZaX}
\\
In this chapter, we designed an online collaborative document editing service. In our design, we provided features like collaborative editing, keeping version history for reverting to older versions, giving users suggestions on frequently used terms and phrases, and view counts of a document. We also assessed the possibility of adding a chatting feature between users collaborating on the same document. A unique aspect of the design was the conflict resolution between concurrent editing operations by different users. We solved the concurrency issues through OT and CRDT.

% End of section content